\section{Results}
\label{sec:results}

All numbers are LLM-as-judge accuracy (qwen2.5-14b judge, temperature 0) over the
full LoCoMo benchmark (1{,}986 QA pairs). Adversarial (category~5) questions are
scored by abstention detection; because a memory-less model abstains on
everything, that category lifts every strategy's floor, so we report it
separately and never fold it into a single headline number. The answerer model is
held fixed (Gemma4-12B-QAT, Q4, on a single 12\,GB RTX~4070).

\subsection{Memory strategies}

\begin{table}[t]
\centering\small
\begin{tabular}{lcccccc}
\toprule
strategy & overall & multi-hop & temporal & open & single-hop & abstain \\
\midrule
none & 0.228 & 0.00 & 0.01 & 0.00 & 0.00 & 1.00 \\
raw\_rag & 0.434 & 0.18 & 0.29 & 0.10 & 0.39 & 0.85 \\
full\_history & \textbf{0.544} & \textbf{0.39} & 0.24 & \textbf{0.22} & \textbf{0.66} & 0.72 \\
agentic (ours) & 0.457 & 0.31 & \textbf{0.37} & 0.16 & 0.35 & 0.88 \\
\bottomrule
\end{tabular}
\caption{Judged accuracy by memory strategy. Agentic memory gives the best
temporal reasoning; full-history wins overall on in-window single-hop copy.}
\label{tab:strategies}
\end{table}

Agentic memory gives the best temporal reasoning of any strategy (0.37 vs.\ 0.24
for full\_history), by grounding absolute dates at write time
(\S\ref{sec:results-temporal}). full\_history wins overall accuracy, driven by
single-hop (0.66) and multi-hop (0.39) questions whose answer is present verbatim
in the window---an advantage that is contingent on the answer surviving
truncation (\S\ref{sec:results-length}).

\subsection{Temporal grounding ablation}
\label{sec:results-temporal}
Absolutizing relative dates at extraction and stamping memories with session
dates is a single isolable mechanism; disabling it drops agentic temporal
accuracy to near zero, enabling it yields the 0.37 above.
\todo{Figure 3: temporal accuracy, grounding off vs.\ on --- from the ablation run.}

\subsection{Length-invariance: the core result}
\label{sec:results-length}

\begin{figure}[t]
\centering
\includegraphics[width=0.8\linewidth]{fig1_length_invariance.pdf}
\caption{Retrieval-based memory is length-invariant; context-stuffing collapses
past its window. full\_history's accuracy drops 76\% once a question's evidence
is truncated, while agentic memory is unchanged.}
\label{fig:length}
\end{figure}

full\_history retains only the most recent 60k characters. Splitting the
questions (excluding adversarial) by whether their evidence still lies in that
window (Table~\ref{tab:length}):

\begin{table}[t]
\centering\small
\begin{tabular}{lcc}
\toprule
strategy & evidence in-window ($n{=}997$) & evidence truncated ($n{=}539$) \\
\midrule
full\_history & 0.671 & 0.163 (\textbf{$-$76\%}) \\
agentic (ours) & 0.337 & 0.330 (\textbf{flat}) \\
raw\_rag & 0.305 & 0.334 (flat) \\
\bottomrule
\end{tabular}
\caption{Judged accuracy by whether evidence survives the context window.}
\label{tab:length}
\end{table}

full\_history's entire advantage comes from questions still inside its window;
beyond it, accuracy collapses 76\%---below either retrieval method. Agentic memory
is length-invariant (0.337 vs.\ 0.330), and outperforms context-stuffing
2.0$\times$ on truncated-evidence questions (Figure~\ref{fig:length}). As a
conversation grows, the truncated fraction approaches 100\%: the regime where
stuffing wins is the one a persistent companion grows out of.

\subsection{Context efficiency}
Agentic memory feeds 889 characters of memory per question vs.\ 59{,}585 for
full\_history---\textbf{67$\times$ less context} for competitive accuracy,
fitting a 4k-token window on a 12\,GB GPU.

\subsection{How small can the memory manager be?}
\label{sec:results-ladder}

\begin{table}[t]
\centering\small
\begin{tabular}{lccccc}
\toprule
extractor & overall & multi-hop & temporal & open & single-hop \\
\midrule
qwen2.5:3b & 0.309 & 0.12 & 0.06 & 0.05 & 0.16 \\
qwen2.5:7b & 0.348 & 0.15 & 0.19 & 0.12 & 0.21 \\
qwen2.5:14b & 0.470 & 0.29 & 0.34 & 0.12 & 0.40 \\
\bottomrule
\end{tabular}
\caption{Judged accuracy by extractor size (answerer fixed). Quality scales with
size; temporal grounding breaks first under compression.}
\label{tab:ladder}
\end{table}

Varying the extractor (answerer fixed), quality scales monotonically with size,
and temporal grounding is the first capability to break: temporal accuracy is
0.06 at 3B, 0.19 at 7B, 0.34 at 14B. A 3B extractor cannot reliably run the typed
protocol on dated information; competence emerges at $\sim$7--14B, setting a
concrete minimum viable size for the memory manager.

\subsection{Extraction precision}
A hand-verified sample of 200 stored memories, traced to source turns across all
conversations, yields \textbf{199/200 faithfulness precision}: zero
hallucinations and one speaker misattribution (a fact-dense speaker's basketball
history duplicated onto another participant). A stricter usefulness bar lowers
this only to 198/200. We report precision directly and treat downstream QA
accuracy as the recall proxy (missed facts surface as wrong answers); a dedicated
recall annotation with inter-annotator agreement is future work.

\subsection{Cost and interactivity}
Extraction runs asynchronously after the reply streams, so it never delays
time-to-first-token. On the RTX~4070: TTFT $\approx$0.9\,s, generation
$\approx$40+ tok/s, extraction adds $\approx$2\,GB VRAM. The full recall--generate--extract
loop runs in real time on a single consumer GPU.
